\clearpage
\fancyhead[L]{\small GAUSS TO BODE.}
\hypertarget{native-444-part-lower}{}
\begin{center}
{\large GAUSS TO BODE.}

\noindent\rule{\linewidth}{0.4pt}

Astronomisches Jahrbuch for 1825. pp. 103..105. Berlin 1821.

\noindent\rule{\linewidth}{0.4pt}
\end{center}

\begin{flushright}
Göttingen, 26 December 1821.
\end{flushright}

You will excuse my long silence when I tell you that I was away from here for most of the year and even made a trip to Altona in late autumn to take delivery of \textsc{Ramsden}’s zenith sector. I returned from there only a short while ago.

In my triangulation, for which I have so far measured the angles at five vertices, I have tried to make the triangles as large as possible. You will have read in no. 126 of our Göttingische gelehrte Anzeigen [vol.\,IV of the Works.] about the new instrument I have been using for this purpose, the \emph{heliotrope}, and the first large-scale trials made with it. Since then I have used it continually, as a target for measuring angles and with no less success for telegraphic signaling. The sunlight reflected from a mirror 2 inches wide and $1\tfrac{1}{4}$ inches high was seen with the \emph{naked eye} at distances of 5, 6, $7\tfrac{1}{2}$, and even once at $9\tfrac{1}{3}$ geographical miles. Its powerful effect usually astonishes those who see it for the first time and have not been prepared for it by theoretical calculation. With even moderately favorable atmospheric conditions, the only limits now on the lengths of the sides of the triangles are those set by the curvature of the Earth. This is especially true if the mirrors are made somewhat larger, as I have done with two new heliotropes of quite different designs.

The activities I have mentioned have so far kept me from attending to the calculation of an ephemeris for the next appearance of Pallas and from putting the earlier observations I made here in order for your yearbook. For now I can enclose only my observations of this year’s comet.

\Needspace{15\baselineskip}
\begin{center}
\small
\setlength{\tabcolsep}{4pt}
\begin{tabular}{l|rrr|rrr|rrr}
1821 & \multicolumn{3}{c|}{\begin{tabular}{c}Mean time\\ at Göttingen\end{tabular}} & \multicolumn{3}{c|}{Right ascension} & \multicolumn{3}{c}{Northern declination} \\
\hline
Jan. 30 & 7\textsuperscript{u} & 34\textsuperscript{m} & 42\textsuperscript{s} & $359^\circ$ & $27'$ & $7''$ & $16^\circ$ & $4'$ & $36''$ \\
Feb. 3 & 7 & 3 & 56 & 359 & 3 & 54 & 15 & 46 & 3 \\
 7 & 6 & 42 & 38 & 358 & 45 & 5 & 15 & 29 & 49 \\
 9 & 6 & 42 & 30 & 358 & 36 & 24 & 15 & 21 & 22 \\
 10 & 6 & 52 & 27 & 358 & 32 & 19 & 16 & 17 & 50 \\
 11 & 7 & 12 & 0 & 358 & 28 & 27 & 15 & 14 & 26 \\
Mar. 1 & 7 & 18 & 7 & 357 & 13 & 34 & 14 & 8 & 37 \\
 5 & 7 & 5 & 28 & 356 & 54 & 11 & 13 & 43 & 5 \\
\end{tabular}
\end{center}

\Needspace{18\baselineskip}
\SourceLeaf{445}
The parabolic elements that Mr. von \textsc{Staudt} calculated from these observations are as follows.

\begin{center}
\small
\begin{tabular}{p{.43\linewidth}p{.49\linewidth}}
Time of perihelion \mbox{.\,.\,.\,.\,.\,.\,.} & 1821 March $21.61890$ at Göttingen \\
—\quad Arg. of latitude \mbox{.\,.\,.\,.\,.\,.} & $169^\circ\;10'\;9''3$ \\
—\quad Logarithm of distance \mbox{.\,.\,.} & $8.9645990$ \\
{\rightskip=0pt plus 3em Longitude of ascending node \mbox{.\,.\,.\,.\,.}\par} & $48^\circ\;44'\;14''7$ from mean equinox of 1 Feb. \\
Inclination of orbit \mbox{.\,.\,.\,.\,.\,.\,.\,.\,.\,.} & $106\quad40\quad16.4$ \\
\end{tabular}
\end{center}

Mr. von \textsc{Staudt} has compared these elements with my observations and also with all the others that had become known at the time. I enclose the results of this comparison for your yearbook in his own hand, since there is too little time for me to copy them out.

\Needspace{34\baselineskip}
\begin{center}
\begin{tabular}{l|r|r}
& \multicolumn{2}{c}{Differences in} \\
& R.A. & Decl. \\
\multicolumn{3}{c}{Paris} \\
Jan. 21 & $+0''2$ & $+0''1$ \\
\multicolumn{3}{c}{Göttingen} \\
Jan. 30 & $-18''8$ & $+34''9$ \\
Feb. 3 & $+3.9$ & $-6.0$ \\
7 & $-13.1$ & $-36.0\!:$ \\
9 & $-0.1$ & $+14.7$ \\
10 & $+2.6$ & $+7.3$ \\
11 & $-1.2$ & $-3.5$ \\
Mar. 1 & $-0.3$ & $-0.3$ \\
5 & $-5.3$ & $-20.6$ \\
\multicolumn{3}{c}{Bremen} \\
Jan. 30 & $-12.6$ & $+12.4$ \\
30 & $+9.1$ & $+34.2$ \\
Feb. 2 & $+24.7$ & $+9.7$ \\
5 & $-7.3$ & $-42.3\!::$ \\
7 & $+8.8$ & $+16.0\!:$ \\
8 & $+4.9$ & $+22.4$ \\
9 & $+5.0$ & $+14.1$ \\
10 & $-5.8$ & $+20.2$ \\
11 & $+4.4$ & $+3.2$ \\
12 & $-9.5$ & $+0.4$ \\
13 & $+0.6$ & $-24.6$ \\
14 & $-2.8$ & $-19.6$ \\
19 & $+10.2$ & $-3.4$ \\
Mar. 1 & $+8.5$ & $-9.4$ \\
5 & $-0.6$ & $-7.6$ \\
6 & $+8.6$ & $-0.2$ \\
\end{tabular}
\end{center}

\Needspace{38\baselineskip}
\begin{center}
\begin{tabular}{l|r|r}
& \multicolumn{2}{c}{Differences in} \\
& R.A. & Decl. \\
\multicolumn{3}{c}{Mannheim} \\
Feb. 6 & $-0''2$ & $+28''9\!:$ \\
7 & $+2.2$ & $+1.3$ \\
8 & $+9.6$ & $-1.1$ \\
9 & $-6.2$ & $+5.7$ \\
10 & $+5.9$ & $+7.3$ \\
11 & $+12.0$ & $+17.4$ \\
12 & $+8.6$ & $+6.3$ \\
13 & $+8.6$ & $+4.2$ \\
14 & $+9.2$ & $+12.7$ \\
15 & $+9.4$ & $-3.7$ \\
27 & $+3.7$ & $-18.2$ \\
\multicolumn{3}{c}{Königsberg} \\
Feb. 9 & $-1.2$ & $+0.5$ \\
10 & $-5.4$ & $+2.9$ \\
11 & $+1.2$ & $-7.6$ \\
12 & $+1.9$ & $-1.1$ \\
14 & $+14.3$ & $-6.3$ \\
15 & $+11.7$ & $-20.8$ \\
19 & $-9.5$ & $+20.7$ \\
25 & $+35.3$ & $+35.9$ \\
27 & $+5.1$ & $-10.6$ \\
Mar. 4 & $+15.6$ & $-12.4$ \\
5 & $-11.4$ & $+4.3$ \\
6 & $+7.8$ & $-17.2$ \\
\multicolumn{3}{c}{Hamburg} \\
Feb. 7 & $-3.2$ & $-29.2$ \\
8 & $+4.5$ & $+28.9$ \\
9 & $-6.5$ & $+9.9$ \\
10 & $+4.5$ & \mbox{.\,.\,.\,.} \\
12 & $-7.0$ & \mbox{.\,.\,.\,.} \\
\multicolumn{3}{c}{Dorpat} \\
Feb. 14 & $-3.6$ & $+83.0\!::$ \\
\end{tabular}
\end{center}

\Needspace{39\baselineskip}
\begin{center}
\begin{tabular}{l|r|r}
& \multicolumn{2}{c}{Differences in} \\
& R.A. & Decl. \\
\multicolumn{3}{c}{Seeberg} \\
Feb. 3 & $+30''2$ & $+103''4\!::$ \\
5 & $-8.6$ & $-41.5$ \\
7 & $+14.3$ & $-95.1\!::$ \\
8 & $+25.1$ & $+40.3\!::$ \\
9 & $+9.8$ & $-20.6$ \\
10 & $-4.9$ & $-19.0$ \\
11 & $+0.1$ & $-15.1$ \\
12 & $+10.2$ & $-20.0$ \\
14 & $+5.3$ & $-39.7$ \\
\multicolumn{3}{c}{Vienna} \\
Feb. 20 & $-78.8$ & $-28.1$ \\
\multicolumn{3}{c}{Milan} \\
Jan. 31 & $-37.7$ & $+4.5$ \\
Feb. 1 & $-22.9$ & $-4.4$ \\
2 & $-0.2$ & $-3.1$ \\
3 & $-29.4$ & $-0.2$ \\
7 & $-31.5$ & $-10.5$ \\
8 & $-19.8$ & $-2.3$ \\
\multicolumn{3}{c}{Padua} \\
Feb. 16 & $-50.5$ & $-56.1$ \\
16 & $-30.8$ & $-60.2$ \\
17 & $-69.4$ & $-15.5$ \\
17 & $-44.0$ & $-45.9$ \\
18 & $-61.5$ & $-1.3$ \\
18 & $-34.6$ & $-30.1$ \\
\multicolumn{3}{c}{Florence} \\
Feb. 16 & $-14.8$ & $-24.9$ \\
17 & $-7.0$ & $-12.9$ \\
21 & $+5.9$ & $-47.6$ \\
22 & $-19.1$ & $-34.8$ \\
23 & $-8.6$ & $-47.5$ \\
25 & $-26.3$ & $-265.0$ \\
27 & $+16.0$ & $-603.2$ \\
\end{tabular}
\end{center}
